\section{Related Work}
\label{sec:related_work}

\paragraph{Temporal alignment for event retrieval.}
Ordered multi-event search has been a central theme of recent Ho Chi Minh City AI Challenge systems.
DANTE builds an event--keyframe similarity matrix and finds a monotone path per video with a gap penalty and a running maximum~\cite{luu2025dante}; TARS uses a prefix-maximum DP recurrence~\cite{do2025tars}; KPTER re-ranks with K pointers~\cite{dinh2025kpter}; the Cross Segment Coherence Scorer adds soft penalties for order reversals, long gaps, and shot jumps~\cite{nguyen2025cscs}; and agent-guided DP lets an LLM coordinate the optimization~\cite{nguyen2025agentdp}.
MADTempo combines query decomposition with beam search~\cite{vu2025madtempo}, and Lucifer-TRACE complements DP with LVLM verification that explicitly does not alter retrieval scores~\cite{ho2025lucifertrace}.
Outside this venue, training-free temporal grounding decomposes a query into sub-events and filters VLM proposals by their order~\cite{zheng2024tfvtg}.
All of these works report end-task scores.
We use the same family of pipeline and instead measure per-event correctness and the effect of each layer.

\paragraph{Query performance prediction.}
Predicting whether a query will succeed is well studied in text retrieval, from query clarity~\cite{cronen2002predicting} to post-retrieval predictors based on the score distribution, such as weighted information gain~\cite{zhou2007wig} and normalized query commitment~\cite{shtok2012nqc}; see~\cite{carmel2010difficulty} for a survey.
For video search, Kofler et al.\ predicted query failure from search logs and visual features of the result list, with failure defined by the absence of clicks~\cite{kofler2012failure}; VQPP introduced a benchmark of single-sentence text-to-video queries and also measured query reformulation~\cite{lutu2026vqpp}.
Unlike these settings, each ranked item here is a jointly decoded multi-event path, the target is whether the top video is correct, and no session history is available; we test how far the simplest post-retrieval signals carry over.

\paragraph{Interaction and correction.}
Interactive video search systems refine their results through relevance feedback, conversation, and query history~\cite{khan2024exquisitor,sharma2025exquisitor,ma2022ivcmr}.
At SoICT 2025, interactive mechanisms refine the query: an LLM agent through dialogue~\cite{mai2025conversational}, composed-image textual feedback~\cite{to2025cirfeedback}, or Rocchio feedback~\cite{nguyen2025vortex}.
Our correction acts on the decoded path rather than on the query.
Constrained decoding with simulated user corrections, propagation to uncorrected fields, and confidence-ordered inspection was introduced for interactive information extraction~\cite{kristjansson2004constrained}; we apply the same idea to ordered video events and measure which of these components matter there.
Building on the EventTrail component of our VBS system SHI~\cite{nguyen2027shi}, we isolate the contribution of the DP starting path, constrained re-decoding, proposal generation, inspection order, and propagation.
